\documentclass[10pt,a4paper]{article}
\usepackage[utf8]{inputenc}
\usepackage[vietnamese]{babel}
\usepackage{amsmath,amssymb,amsfonts}
\usepackage{geometry}
\geometry{a4paper, left=10mm, right=10mm, top=10mm, bottom=10mm}
\usepackage{tikz}
\usetikzlibrary{shapes,arrows.meta,positioning,calc}
\usepackage[most]{tcolorbox}
\usepackage{array,booktabs,multicol}
\usepackage{fontawesome}
\usepackage{helvet}
\renewcommand{\familydefault}{\sfdefault}
\usepackage{parskip}

\definecolor{primary}{RGB}{112, 48, 160}
\definecolor{teal}{RGB}{0, 128, 128}
\definecolor{orange}{RGB}{230, 108, 44}
\definecolor{green}{RGB}{46, 139, 87}
\definecolor{rose}{RGB}{192, 0, 0}

\pagestyle{empty}

\begin{document}

% --- HEADER BANNER ---
\begin{tcolorbox}[
  enhanced,
  colback=violet!8!white,
  colframe=violet!80!black,
  arc=2.5mm, boxrule=1pt,
  left=3mm, right=3mm, top=2mm, bottom=2mm
]
  \begin{tabular*}{\textwidth}{@{\extracolsep{\fill}} lr}
    {\small\bfseries\color{violet!85!black}\faBook\ SỔ TAY KIẾN THỨC TOÁN 2} & 
    {\small\bfseries\color{teal!85!black}\faGraduationCap\ CHƯƠNG TRÌNH GDPT 2018}
  \end{tabular*}
  \vspace{-1mm}
  \begin{center}
    \rule{0.98\textwidth}{0.6pt}
  \end{center}
  \vspace{-2mm}
  \centering
  {\LARGE\bfseries\color{violet!90!black}\faBookmark\ BÀI 52: VIẾT SỐ THÀNH TỔNG CÁC TRĂM, CHỤC, ĐƠN VỊ}\\[1mm]
  {\itshape\footnotesize\color{black!75} Kĩ thuật phân tích số và viết số thành tổng các giá trị theo hàng}
\end{tcolorbox}

\vspace{-2mm}

\noindent
\begin{minipage}[t]{0.49\textwidth}\vspace{0pt}
\begin{tcolorbox}[
  enhanced,
  colback=violet!4!white,
  colframe=violet!75!black,
  arc=2mm, boxrule=0.8pt,
  left=2mm, right=2mm, top=1.8mm, bottom=1.8mm,
  title={\footnotesize\bfseries\color{white}\faKey\ 1. Giá trị theo vị trí của từng chữ số}
]
{\footnotesize
\begin{itemize}\setlength{\itemsep}{1pt}\setlength{\leftskip}{-2.5mm}
  \item Trong số có ba chữ số $\overline{abc}$:
  \begin{itemize}\setlength{\itemsep}{0.5pt}\setlength{\leftskip}{-2mm}
    \item Chữ số $a$ ở hàng trăm có giá trị là $a \times 100$ ($a$ trăm).
    \item Chữ số $b$ ở hàng chục có giá trị là $b \times 10$ ($b$ chục).
    \item Chữ số $c$ ở hàng đơn vị có giá trị là $c$ ($c$ đơn vị).
  \end{itemize}
  \item \textbf{Công thức tổng quát:}
  $$\mathbf{\overline{abc} = a00 + b0 + c}$$
  \item \textbf{Ví dụ:} Số $358$ gồm $3$ trăm, $5$ chục và $8$ đơn vị:
  $$358 = 300 + 50 + 8$$
\end{itemize}
}
\end{tcolorbox}
\end{minipage}
\hfill
\begin{minipage}[t]{0.49\textwidth}\vspace{0pt}
\begin{tcolorbox}[
  enhanced,
  colback=teal!4!white,
  colframe=teal!75!black,
  arc=2mm, boxrule=0.8pt,
  left=2mm, right=2mm, top=1.8mm, bottom=1.8mm,
  title={\footnotesize\bfseries\color{white}\faSliders\ 2. Các trường hợp có chữ số 0}
]
{\footnotesize
\begin{itemize}\setlength{\itemsep}{1pt}\setlength{\leftskip}{-2.5mm}
  \item \textbf{Hàng chục bằng 0:}
  $$\overline{a0c} = a00 + c$$
  \textit{Ví dụ:} $407 = 400 + 7$ (không cần viết $+ 0$).
  \item \textbf{Hàng đơn vị bằng 0 (số tròn chục):}
  $$\overline{ab0} = a00 + b0$$
  \textit{Ví dụ:} $560 = 500 + 60$.
  \item \textbf{Ngược lại: Ghép tổng thành số:}
  $$700 + 30 + 4 = \mathbf{734}$$
  $$900 + 5 = \mathbf{905}\quad (\text{hàng chục là } 0)$$
\end{itemize}
}
\end{tcolorbox}
\end{minipage}

\vspace{-1.5mm}
\begin{tcolorbox}[
  enhanced,
  colback=orange!4!white,
  colframe=orange!80!black,
  arc=2mm, boxrule=0.8pt,
  left=2mm, right=2mm, top=1.5mm, bottom=1.5mm,
  title={\footnotesize\bfseries\color{white}\faSitemap\ 3. Sơ đồ tư duy phân tích cấu tạo số thành tổng}
]
\centering
\begin{tikzpicture}[x=1.5cm, y=0.8cm]
  \tikzset{arr/.style={-{Stealth[scale=0.8]}}, every node/.append style={align=center}}
  \node[rectangle, rounded corners=2pt, draw=violet!80!black, fill=violet!20, font=\footnotesize\bfseries] (num) at (0,0) {Số $\mathbf{642}$};
  \node[rectangle, rounded corners=2pt, draw=teal!80!black, fill=teal!15, font=\scriptsize\bfseries] (tram) at (2.2,0.8) {Hàng trăm: $6$\\Giá trị: $\mathbf{600}$};
  \node[rectangle, rounded corners=2pt, draw=orange!80!black, fill=orange!15, font=\scriptsize\bfseries] (chuc) at (2.2,0) {Hàng chục: $4$\\Giá trị: $\mathbf{40}$};
  \node[rectangle, rounded corners=2pt, draw=green!70!black, fill=green!15, font=\scriptsize\bfseries] (dv) at (2.2,-0.8) {Hàng đ.vị: $2$\\Giá trị: $\mathbf{2}$};

  \node[rectangle, rounded corners=2pt, draw=red!75!black, fill=red!15, font=\footnotesize\bfseries] (res) at (4.5,0) {
    $\mathbf{642 = 600 + 40 + 2}$
  };

  \draw[arr, thick] (num) |- (tram);
  \draw[arr, thick] (num) -- (chuc);
  \draw[arr, thick] (num) |- (dv);
  \draw[arr, thick] (tram) -| (res);
  \draw[arr, thick] (chuc) -- (res);
  \draw[arr, thick] (dv) -| (res);
\end{tikzpicture}
\end{tcolorbox}

\vspace{-1.5mm}
\noindent
\begin{minipage}[t]{0.49\textwidth}\vspace{0pt}
\begin{tcolorbox}[
  enhanced,
  colback=green!4!white,
  colframe=green!60!black,
  arc=2mm, boxrule=0.8pt,
  left=2mm, right=2mm, top=1.8mm, bottom=1.8mm,
  title={\footnotesize\bfseries\color{white}\faCheckCircle\ 4. Ví dụ mẫu vận dụng}
]
{\footnotesize
\textbf{1. Viết các số thành tổng trăm, chục, đơn vị:}\\
- $475 = 400 + 70 + 5$.\\
- $820 = 800 + 20$.\\
- $309 = 300 + 9$.\\
\textbf{2. Viết tổng thành số có 3 chữ số:}\\
- $200 + 60 + 3 = \mathbf{263}$.\\
- $500 + 80 = \mathbf{580}$.\\
- $700 + 4 = \mathbf{704}$ (chú ý hàng chục là $0$).
}
\end{tcolorbox}
\end{minipage}
\hfill
\begin{minipage}[t]{0.49\textwidth}\vspace{0pt}
\begin{tcolorbox}[
  enhanced,
  colback=rose!4!white,
  colframe=rose!70!black,
  arc=2mm, boxrule=0.8pt,
  left=2mm, right=2mm, top=1.8mm, bottom=1.8mm,
  title={\footnotesize\bfseries\color{white}\faExclamationTriangle\ 5. Lỗi sai học sinh hay mắc}
]
{\footnotesize
\begin{itemize}\setlength{\itemsep}{1.5pt}\setlength{\leftskip}{-2.5mm}
  \item \textbf{Quên số 0 ở hàng chục:} Khi viết $700 + 4$, học sinh thường viết nhầm thành $74$ thay vì đúng là $\mathbf{704}$.
  \item \textbf{Nhầm lẫn giữa chữ số và giá trị:}\\
  Trong số $358$, chữ số hàng trăm là $3$, nhưng giá trị của nó là $\mathbf{300}$.
  \item Luôn kiểm tra lại bằng cách cộng nhẩm các số hạng xem có đúng bằng số ban đầu không.
\end{itemize}
}
\end{tcolorbox}
\end{minipage}


\end{document}
